\documentclass{ppfcmthesis}

% Polish captions for contents and lists
\addto\captionspolish{%
  \renewcommand{\contentsname}{Spis treści}%
  \renewcommand{\listfigurename}{Spis rysunków}%
  \renewcommand{\listtablename}{Spis tabel}%
  \renewcommand{\lstlistingname}{Listing}%
  \renewcommand{\lstlistlistingname}{Spis listingów}%
}

% Ensure Polish is active (defensive)
\selectlanguage{polish}

% Thesis information
\thesistitle{Analiza bezpieczeństwa aplikacji mobilnych Android}
\thesisauthor{Mikołaj Kmieciak}
\thesisyear{2025}

\begin{document}

% Title page
\maketitle

% Table of contents and lists (now with consistent arabic numbering)
\tableofcontents
\newpage
\listoffigures
\newpage
\listoftables
\newpage
\lstlistoflistings

\chapter{Wprowadzenie}

To jest pierwszy rozdział pracy magisterskiej. Tutaj przedstawiamy wprowadzenie do tematu.

\section{Cele pracy}

Główne cele tej pracy to:
\begin{itemize}
\item Analiza bezpieczeństwa aplikacji Android
\item Opracowanie metodologii testowania
\item Implementacja narzędzi testowych
\end{itemize}

\begin{figure}[h]
\centering
\rule{5cm}{3cm}
\caption{Przykładowy rysunek}
\label{fig:przyklad}
\end{figure}

\begin{table}[h]
\centering
\begin{tabular}{|c|c|}
\hline
Kolumna 1 & Kolumna 2 \\
\hline
Wartość 1 & Wartość 2 \\
\hline
\end{tabular}
\caption{Przykładowa tabela}
\label{tab:przyklad}
\end{table}

\begin{lstlisting}[language=Java, caption=Przykładowy kod Java]
public class Example {
    public static void main(String[] args) {
        System.out.println("Hello World!");
    }
}
\end{lstlisting}

\begin{equation}
E = mc^2
\label{eq:einstein}
\end{equation}

\chapter{Przegląd literatury}

Drugi rozdział zawiera przegląd literatury przedmiotu.

\section{Stan badań}

Obecny stan badań w dziedzinie bezpieczeństwa Android...

\chapter{Metodologia}

Trzeci rozdział opisuje metodologię badawczą.

\chapter{Implementacja}

Czwarty rozdział przedstawia szczegóły implementacji.

\chapter{Wyniki}

Piąty rozdział zawiera wyniki przeprowadzonych badań.

\chapter{Wnioski}

Ostatni rozdział podsumowuje pracę i przedstawia wnioski.

\end{document}